\gdef\tagBrand{NT}
\gdef\footerAddress{Zalo: 0799666070}
\setcounter{ex}{0}

% --- CÂU 1 ---
\begin{ex}
	Viết bất phương trình $-4x + 6y \le 18$ theo dạng $ax + by \le c$. Giá trị của $c$ bằng bao nhiêu?
	\shortans{18}
	\loigiai{
		So sánh bất phương trình $-4x + 6y \le 18$ với dạng tổng quát $ax + by \le c$, ta có $a = -4$, $b = 6$ và $c = 18$.\\
		Vậy giá trị của $c$ bằng $18$.
	}
\end{ex}

% --- CÂU 2 ---
\begin{ex}
	Hệ bất phương trình gồm $x \ge 0$, $y \ge 0$, $3x + 2y \le 12$, $x + 4y \le 10$ có bao nhiêu bất phương trình thành phần?
	\shortans{4}
	\loigiai{
		Hệ bất phương trình đã cho bao gồm $4$ bất phương trình thành phần:
		\begin{enumerate}[(1)]
			\item $x \ge 0$;
			\item $y \ge 0$;
			\item $3x + 2y \le 12$;
			\item $x + 4y \le 10$.
		\end{enumerate}
		Vậy hệ có $4$ bất phương trình thành phần.
	}
\end{ex}

% --- CÂU 3 ---
\begin{ex}
	Tính giá trị của biểu thức $M = \sin 30^\circ + \cos 60^\circ$.
	\shortans{1}
	\loigiai{
		Ta có:
		$$\sin 30^\circ = \dfrac{1}{2}\quad \text{và}\quad \cos 60^\circ = \dfrac{1}{2}.$$
		Do đó:
		$$M = \sin 30^\circ + \cos 60^\circ = \dfrac{1}{2} + \dfrac{1}{2} = 1.$$
	}
\end{ex}

% --- CÂU 4 ---
\begin{ex}
	Cho tam giác $ABC$ có cạnh $BC = 12\text{ cm}$ và $\widehat{A} = 30^\circ$. Bán kính đường tròn ngoại tiếp $R$ của tam giác $ABC$ bằng bao nhiêu xentimét?
	\shortans{12}
	\loigiai{
		Áp dụng định lí sin trong tam giác $ABC$, ta có:
		$$\dfrac{BC}{\sin A} = 2R \implies R = \dfrac{BC}{2\sin A} = \dfrac{12}{2\sin 30^\circ} = \dfrac{12}{2 \cdot \dfrac{1}{2}} = 12\text{ (cm)}.$$
		Vậy bán kính đường tròn ngoại tiếp tam giác là $R = 12\text{ cm}$.
	}
\end{ex}
